\documentclass[11pt,a4paper]{article}
\usepackage{amsmath,amssymb,amsthm}
\usepackage[margin=2.5cm]{geometry}
\usepackage{hyperref}

\newtheorem{definition}{Definition}[section]
\newtheorem{theorem}[definition]{Theorem}
\newtheorem{proposition}[definition]{Proposition}
\newtheorem{lemma}[definition]{Lemma}
\newtheorem{corollary}[definition]{Corollary}
\newtheorem{conjecture}[definition]{Conjecture}
\newtheorem{remark}[definition]{Remark}

\title{Hyperseed Formalization:\\
  Evidence Is to Logic What Energy Is to Physics}
\author{ProtomegaTron (automated formalization)}
\date{2026-09-18}

\begin{document}
\maketitle

\begin{abstract}
We formalize the intellectual structure of Ben Goertzel's Substack article
``Evidence Is to Logic What Energy Is to Physics'' (2026-03-10) within the
Hyperseed ontology. The article develops the evidence/energy analogy into
rigorous mathematics: a discrete quantale Noether theorem, five
anti-hallucination theorems, a non-commutative extension yielding quantum
logic, and QLN generalizing PLN. Key mappings: evidence conservation as
fiber conservation, the Noether theorem as fiber-base symmetry,
hallucination bound as fiber-base constraint, non-commutativity as fiber
operation order, and QLN as quantum fiber inference.
\end{abstract}

\tableofcontents

%% =================================================================
\section{Fiber conservation: the quantale Noether theorem}
\label{sec:noether}
%% =================================================================

\begin{theorem}[Evidence conservation --- claims 1--8, 25--26]
\label{thm:noether}
Evidence is to logic what energy is to physics. An inference engine
traversing a proof graph converts premise-evidence into conclusion-evidence,
just as a ball rolling down a hill converts potential energy into kinetic
energy. If the engine is honest --- neither fabricating evidence nor
double-counting it --- then evidence is conserved along optimal inference
paths.

The mathematical setting is quantales: complete lattices equipped with a
monoidal product $\otimes$ that distributes over joins. Different quantale
choices give different logics:
\begin{itemize}
\item Nonnegative reals $(\mathbb{R}_{\geq 0}, +, \times)$: probabilistic reasoning.
\item Tropical semiring $(\mathbb{R} \cup \{\infty\}, \min, +)$: shortest-path problems.
\item Boolean $(\{0,1\}, \lor, \land)$: classical logic.
\end{itemize}

Along an inference path, define forward factor $f$ and backward factor $g$.
Their product $\rho = f \otimes g$ is the ``reinforcement'' at each node.

\textbf{Discrete Quantale Noether Theorem:} Along geodesic (optimal)
inference paths, the reinforcement $\rho$ is constant:
\[
  \rho(n_i) = f_i \otimes g_i = \text{const} \quad
  \text{for all nodes } n_i \text{ on the geodesic}
\]

The proof is purely lattice-theoretic: it uses only that $\otimes$
distributes over joins and that paths through a given node are a subset
of all paths. This is the exact analogue of energy conservation ---
not metaphor but isomorphic mathematical structure.

In Hyperseed terms, this is \emph{fiber conservation}: evidence (fiber
content) is conserved along optimal inference paths (base geodesics).
The fiber can't create evidence from nothing --- it can only transport,
transform, and distribute the evidence that enters from the base:
\[
  \nabla_{\gamma} \rho = 0 \quad \text{along geodesic } \gamma
\]

The conservation follows from \emph{fiber-base symmetry}: invariance
under translation along the inference path. The fiber-base connection
has a symmetry (path-translation invariance), and evidence conservation
is the corresponding conserved quantity --- the Noether correspondence
applied to the fiber bundle.
\end{theorem}

%% =================================================================
\section{Anti-hallucination as fiber-base constraints}
\label{sec:hallucination}
%% =================================================================

\begin{proposition}[Five theorems --- claims 9--13, 27, 32]
\label{prop:hallucination}
Evidence conservation alone doesn't fully constrain inference, just as
energy conservation alone doesn't fully constrain physics. Five additional
theorems complete the picture:

\textbf{1. Hallucination Bound (fiber can't exceed base evidence):}
Conclusion strength is bounded by premise evidence mass. No inference
chain can manufacture support that wasn't there at the start:
\[
  \text{evidence}(\text{conclusion}) \leq \text{evidence}(\text{premises})
\]

In Hyperseed terms: the fiber can't contain more evidence than the base
provides. This is a fundamental fiber-base constraint: no matter how
sophisticated the fiber operations (inference steps), the total evidence
is bounded by the base input.

\textbf{2. Weakness-Bounded Leakage (fiber operation reordering cost):}
Reordering near-commutative inference steps introduces discrepancy
bounded by their pairwise ``weakness'' --- a quantale-valued measure of
how much order matters:
\[
  \|f \circ g - g \circ f\| \leq w(f, g) \quad \text{(weakness bound)}
\]

This makes parallel and asynchronous inference safe: you can schedule
steps freely, and the evidence integrity cost is controlled.

\textbf{3. Evidence Monotonicity (fiber data-processing inequality):}
Capsule-respecting inference cannot increase total evidence content ---
the quantale analog of the classical data-processing inequality.

\textbf{4. Join-Collision Entropy Non-Decrease (fiber second law):}
Under join-bistochastic inference maps, evidence concentration can only
decrease --- the logical second law of thermodynamics:
\[
  H_{\text{join}}(\rho_{t+1}) \geq H_{\text{join}}(\rho_t)
\]

In Hyperseed terms: without new evidence input from the base, the
fiber's evidence becomes increasingly diffuse. The fiber tends toward
maximum uncertainty (maximum entropy) without external input.

\textbf{5. Together:} An honest inference engine conserves evidence along
optimal paths, cannot hallucinate beyond its premises, degrades gracefully
under reordering, and tends toward increasing uncertainty without new input.
These are mathematical consequences of evidence conservation, not
engineering desiderata.
\end{proposition}

%% =================================================================
\section{Fiber operation order: the non-commutative extension}
\label{sec:noncommutative}
%% =================================================================

\begin{theorem}[Non-commutative evidence --- claims 14--17, 28--30]
\label{thm:noncommutative}
When $a \otimes b \neq b \otimes a$ (the quantale is non-commutative),
the order of inference steps matters --- just as the order of measurements
matters in quantum mechanics.

In the non-commutative extension, the Noether theorem splits into left
and right conservation laws:
\[
  \rho_L = f \otimes g = \text{const} \qquad
  \rho_R = g \otimes f = \text{const}
\]

These correspond to left and right fiber currents --- two independent
conservation laws arising from the non-commutativity. In the commutative
case they coincide ($\rho_L = \rho_R$); in the non-commutative case they
diverge.

This yields an \emph{uncertainty principle for inference}: you cannot
simultaneously know the evidence for two non-commuting propositions to
arbitrary precision. In Hyperseed terms, this is fiber measurement
incompatibility: the fiber can't simultaneously resolve two non-commuting
base properties. This connects directly to Level~2 (incompatibility) in
the observer-relative quantumity hierarchy.

The non-commutative extension leads to a new derivation of quantum
mechanics itself: QM emerges from evidence conservation in
non-commutative settings. The Hilbert space structure, the Born rule,
the unitary evolution --- all follow from evidence conservation in a
non-commutative quantale. This is not QM-by-analogy but QM-from-evidence:
the mathematical structure of quantum mechanics is the inevitable
consequence of conserving evidence when operations don't commute.
\end{theorem}

%% =================================================================
\section{Quantum fiber inference: QLN}
\label{sec:qln}
%% =================================================================

\begin{definition}[Quantum Logic Networks --- claims 18--24, 31]
\label{def:qln}
QLN (Quantum Logic Networks) generalizes PLN (Probabilistic Logic
Networks) into the quantum regime:
\begin{itemize}
\item \textbf{Evidence states} are represented as density matrices ---
  mixed states capturing uncertainty about the evidence configuration.
\item \textbf{Inference steps} are completely positive, trace-preserving
  (CPTP) maps --- the quantum analog of stochastic maps.
\item \textbf{Algebraic setting} is a non-commutative quantale --- the
  generalization of the probabilistic quantale that PLN uses.
\end{itemize}

In Hyperseed terms, QLN is the inference engine for \emph{quantum fibers}:
fibers whose operations are inherently non-commutative. PLN is the
commutative special case:
\[
  \text{QLN} \xrightarrow{[\cdot, \cdot] = 0} \text{PLN}
\]

The unification is precise: PLN and QLN share the same evidence
conservation principle (quantale Noether theorem), the same
anti-hallucination guarantees, and the same fiber-conservation
framework. The only difference is whether the fiber operations commute.

QLN + conserved fluidic routing = FluQNets. The evidence/energy analogy
connects to the fluidic neural network architecture through the shared
conservation principle:
\[
  \underbrace{\text{QLN}}_{\text{fiber inference}} +
  \underbrace{\text{conserved routing}}_{\text{base dynamics}} =
  \underbrace{\text{FluQNets}}_{\text{fiber-conserved bundle}}
\]

The HJB-Navier-Stokes-Schr\"odinger mapping chain provides the
mathematical bridge: HJB (dynamic programming on the base) maps to
Navier-Stokes (fluidic routing on the base) and to Schr\"odinger
(quantum evolution on the fiber). All three descriptions are equivalent
--- different perspectives on the same fiber-conserved bundle.

The paper series includes results on algorithmic tractability: QLN
reasoning can be made tractable on appropriate quantum hardware,
suggesting that quantum computers may be natural inference engines
for quantum fiber operations.
\end{definition}

%% =================================================================
\section{Cross-references}
%% =================================================================

\begin{remark}[Connections to other formalizations]
\begin{itemize}
\item \textbf{Sometimes a Smaller System Should Model a Bigger System
  as Quantum} (2026-03-18): observer-relative quantumity. The evidence
  conservation framework provides the mathematical foundation for why
  classical systems exhibit quantum-like behavior from limited observers.
\item \textbf{Rethinking the Classical/Quantum Brain} (2026-03-17):
  FluQNets. The evidence conservation principle is the foundation of
  FluQNets, which the brain article applies to neuroscience.
\item \textbf{Google's Quantum Crypto Paper} (2026-04-01): quantum
  hardware. The tractability results for QLN connect to the quantum
  computing capabilities discussed in the crypto article.
\item \textbf{A New Approach to Grand Unified Physics} (2026-05-12):
  physics connections. The evidence/energy analogy connects to the
  broader physics unification program --- evidence conservation as a
  fundamental physical principle.
\item \textbf{Navier-Stokes Blows Up the Internet} (2026-07-01):
  Navier-Stokes. The HJB-NS mapping chain connects evidence conservation
  to fluid dynamics, which connects to the Navier-Stokes work.
\item \textbf{Seeding RSI} (2026-07-28): Hyperon architecture. QLN/FluQNets
  as potential components of Hyperon's inference engine.
\end{itemize}
\end{remark}

\begin{conjecture}[Evidence conservation as foundation of logic]
Conjecture: evidence conservation is the correct foundation for all of
logic --- classical, probabilistic, quantum, and beyond. Every sound
logic can be derived from an appropriate evidence conservation principle
(quantale Noether theorem) applied to an appropriate quantale.

If this conjecture holds, the classification of logics reduces to the
classification of quantales: each quantale determines a logic through
its Noether theorem, and the anti-hallucination theorems hold
universally. The question ``what logic should an AGI use?'' becomes
``what quantale should an AGI's evidence live in?'' --- and the answer
may be ``a non-commutative quantale,'' which gives QLN.

This would mean that quantum mechanics is not just a physical theory
but a logical necessity: any sufficiently rich inference system operating
in a non-commutative evidence space must exhibit quantum-like structure.
QM is not a fact about physics --- it's a fact about evidence conservation.
\end{conjecture}

\end{document}
